\section{Restrykcyjne środowisko}
Celem pracy jest analiza rozwiązań w kontekście ograniczających warunków. Poniższy rozdział omawia zagadnienia związane z tego typu środowiskiem.

\subsection{Rodzaj NAT}
Jest to decydujący czynnik utrudniający nawiązywanie połączeń z urządzeniami spoza sieci lokalnej. W zależności od typu NAT za jakimi znajdują się urządzenia, będzie to proces utrudniony bądź w pełni uniemożliwiony.

\subsubsection{Static NAT}
Przydziela publiczny adres każdemu urządzeniu z puli publicznych adresów. Z tego powodu nie stanowi bariery przy nawiązywaniu bezpośrednich połączeń między urządzeniami. Niewykorzystywany dla rozwiązań konsumenckich.

\subsubsection{Dynamic NAT}
Rozwiązanie rzadziej stosowane działające podobne do statycznego. Jedyna różnica jest taka, że adresy przydzielane są w sposób dynamiczny.

\newpage

\subsubsection{Network Address and Port Translation (NAPT)}
Umożliwia współdzielenie jednego publicznego adresu IP przez wiele urządzeń. Zasada działania tego mechanizmu opiera się na użyciu portów dla rozróżnienia poszczególnych klientów.\\

Każdy socket (kombinacja adresu i portu) w sieci wewnętrznej ma przypisany swój odpowiednik po publicznej stronie routera.\\

Przypisanie to odbywa się automatycznie, na żądanie, jest tymczasowe oraz w większości losowe.\\

Zapytania wychodzące są przekierowywane przez przypisany socket zewnętrzny (router podmienia adres i port źródłowy zapytań).\\

Analogicznie, wszelkie pakiety przychodzące na przypisany socket zewnętrzny są przekierowywane na socket wewnętrzny (router podmienia adres i port docelowy).\\

Przykład tego procesu został przedstawiony na poniższym diagramie:

\begin{figure}[h]
    \centering
    \includegraphics[width=1\textwidth]{nat_diagram}
    \caption{Schemat działania mechanizmu NAPT}
    \label{fig:natdiagram}
\end{figure}

\newpage

Istnieje kilka implementacji NAPT różniących się poziomem restrykcyjności odnośnie filtrowania przychodzących pakietów.\\

\textbf{Full Cone NAT}\\
Najbardziej permisywna implementacja. Nie wprowadza dodatkowych ograniczeń. Mapowanie portów może być predefiniowane lub definiowane w odpowiedzi na wysyłane zapytania.\\

\textbf{Restricted Cone NAT}\\
W tej implementacji pakiety przychodzące są domyślnie odrzucane. Zapytania wychodzące wyzwalają powstanie tymczasowych reguł akceptujących ruch przychodzący z tego samego adresu co adres docelowy zapytania.\\

\textbf{Port Restricted Cone NAT}\\
W tej implementacji pakiety przychodzące są domyślnie odrzucane. Zapytania wychodzące wyzwalają powstanie tymczasowych reguł akceptujących ruch przychodzący z tego samego socketu (adres:port) co socket docelowy zapytania.\\

\textbf{Symmetric NAT}\\
Najbardziej restrykcyjna implementacja. Posiada te same ograniczenia co poprzednia. Modyfikacja polega na zmianie mechanizmu przypisywania portów zewnętrznych.\\

Wcześniejsze implementacje współdzielą wybrany socket zewnętrzny pomiędzy adresami docelowymi. W tej implementacji każdy socket wewnętrzny otrzymuje dedykowany socket zewnętrzny, dla każdego zewnętrznego hosta do którego wysłał zapytanie.\\

Gdy oba urządzenia znajdują się za tego typu NAT-em, nie ma możliwości nawiązania bezpośredniego połączenia pomiędzy nimi.

\subsubsection{Carrier-grade NAT}
Jest to specyficzny przypadek zastosowania NAPT. Odnosi się on do jego wykorzystania w sieci wewnętrznej dostawcy internetu. Użytkownik współdzieli jeden adres z innymi klientami dostawcy. Istnieje wiele implementacji CGNAT. Natomiast główny podział wynika z wykorzystywanego protokołu.\\

Popularny mechanizm DS-Lite zakłada że klient nie otrzymuje żadnego adresu IPv4 a jedynie IPv6. Pakiety IPv4 są enkapsulowane w pakiety IPv6 przy przesyle do routera brzegowego dostawcy.\\

Wart ponownego odnotowania jest fakt, że jest to specyficzny sposób wykorzystania NAPT, aniżeli jego konkretna implementacja. Przykładowo, CGNAT może być w NAT-em symetrycznym, ale równie dobrze może to być Full Cone NAT.

\newpage

\subsection{Metody obejścia NAT}

\subsubsection{STUN (Session Traversal Utilities for NAT)}

Jest to protokół sieciowy pozwalający klientowi znajdującemu się za NAT poznać swój publiczny adres IP wyjściowy i port.\\

Klient z wewnątrz sieci prywatnej wysyła zapytanie do serwera STUN dostępnego publicznie. Serwer STUN analizuje nagłówki pakietu i odsyła do klienta odpowiedź zawierającą jego publiczny adres IP i port.\\

Jest to fundament tworzenia połączeń P2P.

\subsubsection{UDP Hole Punching}

Ta podstawowa technika polegająca na wykorzystaniu mechaniki działania NAT.\\

W uproszczeniu proces ten wygląda następująco:

\begin{enumerate}[leftmargin=*]
    \item Klienci poznają swoje publiczne adresy oraz porty z wykorzystaniem serwera STUN.
    \item Klienci dokonują wymiany tych informacji za pośrednictwem współdzielonego kanału sygnałowego. Przykładem kanału sygnałowego może być publiczny serwer gry lub backend danej aplikacji.
    \item Obaj klienci rozpoczynają cykliczne wysyłanie pakietów na otrzymany adres partnera do momentu nawiązania połączenia. W zależności od warunków sieciowych oraz stopnia synchronizacji, część pakietów może zostać odrzucona do momentu gdy druga strona również nie rozpocznie wysyłania.
    \item Powstają reguły akceptujące ruch przychodzący od strony przeciwnej, istnieje kanał komunikacji P2P.
\end{enumerate}

Do tego wymagana jest stałość mapowania zewnętrznych portów przez NAPT. W innym przypadku port zewnętrzny zidentyfikowany przez serwer STUN nie będzie się pokrywał z tym który zostanie otwarty dla przeciwległego klienta. Z tego powodu jeżeli obaj klienci znajdują się za symetrycznym NAT-em, nawiązanie połączenia P2P w ten sposób jest niemożliwe.\\

\subsubsection{TURN (Traversal Using Relays around NAT)}

Jest to rozszerzenie protokołu STUN. Umożliwia stworzenie serwera funkcjonującego jako przekaźnik ruchu, w pełni omijając w ten sposób ograniczenia stawiane przez NAT. \\

Protokół ten działa w warstwie siódmej modelu OSI, implementowany jest on przez daną aplikację do obsługi pojedynczego socketu.\\

Z uwagi na to że serwer pośredniczy w komunikacji pomiędzy klientami, wprowadza on dodatkowy narzut sieciowy i wymaga dodatkowych zasobów.

\newpage

\subsection{Wysokie opóźnienia i jitter}
VPN nadaje dodatkowe opóźnienie ponad to bazowe, wynikające z sieci użytkownika. W modelu rozproszonym, opóźnienia są zminimalizowane ponieważ poszczególne urządzenia komunikują się bezpośrednio ze sobą. W modelu scentralizowanym, każdy pakiet musi niejako nadrabiać drogę poprzez przejście przez serwer zanim trafi do odbiorcy.\\

Jest to zależne od zastosowania, jednak opóźnienie może mieć na tyle duży wpływ by kategorycznie wykluczyć wykorzystanie sieci wirtualnych w pewnych przypadkach.

\subsection{Utrata pakietów}
W tunelach UDP, utrata pakietów powoduje braki w danych, ale nie zatrzymuje całej transmisji.\\

Natomiast w tunelach TCP może ona spowodować zjawisko TCP Meltdown - zarówno sam VPN, jak i aplikacja wewnątrz tunelu próbują jednocześnie retransmitować zgubione pakiety. Skutkuje to lawinowym wzrostem opóźnień, drastycznym spadkiem przepustowości i częstym zrywaniem sesji.

\subsection{Ograniczenia protokołu UDP}
Znacząco utrudnia zastosowanie większości z dostępnych rozwiązań. Może obejmować kompletną blokadę lub ograniczenie przepustowości względem TCP.\\

WireGuard oraz Nebula nie wspierają TCP w związku z czym ograniczenia nakładane na protokół UDP kompletnie uniemożliwiają lub utrudniają ich wykorzystanie.\\
OpenVPN oraz strongSwan są częściowo odporne, ponieważ przewidują wsparcie dla protokołu TCP.\\

Mimo tego UDP pozostaje preferowanym protokołem, niezależnie od wsparcia z uwagi na negatywne zjawisko TCP Meltdown oraz mniejszy narzut.

\newpage

\subsection{Ograniczenia wydajności}

Kluczowym czynnikiem wpływającym na przepustowość szyfrowanego połączenia są możliwości kryptograficzne procesora. Dla urządzeń ze słabszymi jednostkami, wydajność protokołu ma jeszcze wyższe znaczenie dla osiąganych wyników.\\

Tam gdzie maksymalna przepustowość ograniczona jest innymi czynnikami niż procesor, wydajność dalej ma znaczenie z uwagi na optymalność wykorzystania dostępnych zasobów.\\

W takim przypadku, mimo że osiągana przepustowość jest podobna, interesującym parametrem staje się poziom wykorzystana procesora, ilość zużywanej energii.

\subsection{Ograniczenia przepustowości}

Ograniczenia stawiane przez infrastrukturę, dyktują maksymalne wyniki przepustowości możliwe do osiągnięcia przez protokół.\\

W przypadku gdy ograniczenia te są wysokie, ważniejszymi kryteriami wyboru stają są inne wady oraz zalety protokołu niż surowe wyniki wydajności.